%\documentclass[16pt]{extarticle} 
\documentclass[12pt]{article}
\usepackage[utf8]{inputenc}
\usepackage[T1]{fontenc}
\usepackage[russian]{babel}
\usepackage{geometry}
\usepackage{titlesec}
\usepackage{tocloft}
\usepackage{listings}
\usepackage{attachfile}
\usepackage{dirtree}
\usepackage[dvipsnames]{xcolor}
\makeatletter
\def\AF@Color{0 0 1}  % Синий (RGB)
\makeatother

\lstset{
    language=bash,        % Используем bash, так как он отлично показывает строки с командами
    basicstyle=\ttfamily\small,
    numbers=left,
    numberstyle=\tiny,
    breaklines=true,
    showstringspaces=false,
    frame=single,
    % Добавим немного цветов для наглядности (опционально)
    commentstyle=\color{gray},
    keywordstyle=\color{blue},
    stringstyle=\color{teal}
}

\geometry{a4paper, left=2.5cm, right=2.5cm, top=2.5cm, bottom=2.5cm}

% Отключение нумерации страниц на титуле
\pagenumbering{gobble}

\begin{document}


\begin{titlepage}
    \centering
    \vspace*{\fill}
    
    {\Huge\bfseries Base Github workflows for Ansible}
    
    \vspace{0.5cm}
    
    {\Large Тимофей Семисоха}
    
    \vspace*{\fill}
\end{titlepage}

\newpage
\pagenumbering{arabic}
\setcounter{page}{2}


{\Large\bfseries Base Github workflows for Ansible} \par
В этом небольшом тексте рассмотрим как можно использовать github workflows для автоматического тестирования ansible плейбуков в репозитории с помощью molecule\par\vspace{10pt}
{\Large\bfseries Структура репозитория} \par
Для работы workflows нужно, чтобы репозиторий был построен по следующему принципу:\\
\dirtree{%
.1 .github/workflows.
.2 molecule.yml.
.1 molecule.
.2 ci.
.3 converge.yml.
.3 molecule.yml.
.3 verify.yml.
.2 default.
.3 converge.yml.
.3 molecule.yml.
.3 verify.yml.
}
Суть следующая - у нас есть основные папки, это .github для настроек, о которых мы сообщаем github. И molecule, в которой соблюдается классическая схема организации папок для Molecule.\par\vspace{10pt}
{\Large\bfseries Папки для Molecule}\par
Внутри указаны базовые группы сценариев, это ci и default. Можно сказать, что это сами тесты, которые мы хотим проводить. Если бы мы хотели создать еще один тест, то в той же папке molecule нужно было бы создать папку "onemoretest".\\
В каждом из таких тестов лежат сценарии, по которым мы собираемся проводить тестирование. В структуре указаны минимальные три файла, с которыми Molecule может отработать. Если нужно провести более серьезное тестирование, в папку к имеющимся сценариям можно добавить остальные, приемлемые для Molecule (например prepare.yml).\\
На десктопе при такой структуре файлов default тест запускался бы командой molecule test. Тест ci запускался бы командой molecule test -s ci
\newpage
{\Large\bfseries Workflows}\par
Перейдем к папке .github с действиями для гитхаба. Для тестов с помощью молекулы нам понадобится файл molecule.yml, в котором мы расскажем гитхабу, каких действий от него ожидаем.\\
\begin{lstlisting}
name: Molecule Test

on:
  push:
    branches: [ main, master ]
  pull_request:
    branches: [ main, master ]

jobs:
  molecule:
    runs-on: ubuntu-latest

    steps:
      - name: Checkout code
        uses: actions/checkout@v4

      - name: Set up Python
        uses: actions/setup-python@v5
        with:
          python-version: "3.12"
        
      - name: Install dependencies
        run: |
          python -m pip install --upgrade pip
          pip install molecule molecule-docker ansible-lint
          pip install --upgrade molecule-docker
          

      - name: Run Molecule tests (ci scenario)
        run: molecule test -s ci
        env:
          PY_COLORS: "1"
          ANSIBLE_FORCE_COLOR: "1"
          ANSIBLE_ALLOW_BROKEN_CONDITIONALS: "true"
\end{lstlisting}
После добавления такого файла в указанную директорию, Github начнет тестировать наш репозиторий с помощью Molecule после каждого изменения в основной ветке. Рассмотрим подробнее, из чего состоит этот файл, и чем конкретно будет заниматься Github.\\
\newpage
\begin{lstlisting}
name: Molecule Test

on:
  push:
    branches: [ main, master ]
  pull_request:
    branches: [ main, master ]
\end{lstlisting}\par
Это верхушка нашего файла. Здесь указано название сценария, под ним мы указываем после каких действий мы хотим проводить тестирование. В примере указаны пуши и пул реквесты в мейн.\par\vspace{10pt}

\begin{lstlisting}
jobs:
  molecule:
    runs-on: ubuntu-latest
\end{lstlisting}\par
Здесь указываем работы, которые нужно будет провести, и с помощью каких инструментов. В нашем случае говорим, что хотим контейнер с ubuntu-latest, в котором будет запущена Molecule. (Контейнеры, которые создает Molecule будут созданы внутри этого контейнера)
\begin{lstlisting}
    steps:
      - name: Checkout code
        uses: actions/checkout@v4
\end{lstlisting}\par
Здесь и далее перечисляются шаги, которые будут выполнены для тестирования.\\
actions / checkout@v4 используется для загрузки файлов из репозитория в контейнер в котором будет запускаться Molecule.\par\vspace{10pt}
\begin{lstlisting}
- name: Set up Python
        uses: actions/setup-python@v5
        with:
          python-version: "3.12"
\end{lstlisting}\par
Здесь Github скачивает python версии "3.12" и добавляет его в PATH. Все последующие команды будут выполняться этим python, а не встроенным в Ubuntu.\par\vspace{10pt}
\begin{lstlisting}
      - name: Install dependencies
        run: |
          python -m pip install --upgrade pip
          pip install molecule molecule-docker ansible-lint
          pip install --upgrade molecule-docker
\end{lstlisting}\par
Скачиваем Molecule, драйвер для Docker и утилиту для статического анализа скриптов ansible\par\vspace{10pt}
\newpage

\begin{lstlisting}
      - name: Run Molecule tests (ci scenario)
        run: molecule test -s ci
        env:
          PY_COLORS: "1"
          ANSIBLE_FORCE_COLOR: "1"
          ANSIBLE_ALLOW_BROKEN_CONDITIONALS: "true"
\end{lstlisting}\par
Запускаем тест ci. Ниже устанавливаются цвета для логов. 
Дополнительно здесь указан параметр для замены ошибок, связанных с некорректыми условиями, на предупреждения. По почти подтвержденной информации, в старых версиях скриптов можно было полагаться на истинность небулевых типов, а сейчас так нельзя. Для совместимости со старой логикой можно воспользоваться ANSIBLE BROKEN CONDITIONALS. \\
Возможно, это было важно именно для моего плейбука. Но если столкнетесь с ошибками логики, отдаленно связанными с вашим ansible скриптом, можете попробовать добавить эту строку.\par\vspace{10pt}





\end{document}
